% GENERATED FILE. Edit canonical lessons, then run npm run generate:books.
% canonical-source-hash: fnv1a64:9b7ffa564ca8f956
% canonical-lessons: SA-C73-sings, SA-C73-tells, SA-C73-cries, SA-C73-calls, SA-C73-bows

\chapter{Sing, Tell, Cry, Call}
\label{ch:sa-sing-tell-cry-call}
\hlchaptermodality{73}

\begin{chapteropening}
\textbf{By the end of this chapter:} \emph{I can say someone sings, tells, cries, calls or bows.}

Complete the last lesson of chapter 73: i can say someone sings, tells, cries, calls or bows.
\end{chapteropening}

\section[gāyati]{\sk{गायति} (gāyati) — sings}

\label{lesson:SA-C73-sings}



\begin{quote}
Before the new one: say the Sanskrit for he brings, then the Sanskrit for she gives.
\end{quote}

\subsection*{You'll want to know: \sk{गायति}}
\textbf{\sk{गायति}} (\emph{gāyati}) — \textquotedblleft{}(he, she) sings\textquotedblright{}. What is sung is a \textbf{\sk{गीतम्}}.

\begin{grammarlens}[title={What you've built}]
The verb of a song.
\end{grammarlens}

\subsection*{Your turn}
\begin{itemize}
\raggedright
  \item \emph{Say it:} \emph{gāyati}
  \item \emph{Say it:} \emph{gāyati}, once more
  \item \emph{Say it:} say \emph{gītaṁ gāyati}
  \item \emph{Recall:} say the Sanskrit for he brings, then the Sanskrit for she gives, then say \emph{gāyati} again
\end{itemize}

\begin{tcolorbox}[breakable,colback=teal!4,colframe=teal!35!black,title={Before you move on}]
What does \sk{गायति} mean? (\textquotedblleft{}Sings\textquotedblright{}.) Say it once more.
\end{tcolorbox}
\section[kathayati]{\sk{कथयति} (kathayati) — tells}

\label{lesson:SA-C73-tells}



\begin{quote}
Before the new one: say the Sanskrit for she gives, then the Sanskrit for she sings.
\end{quote}

\subsection*{You'll want to know: \sk{कथयति}}
\textbf{\sk{कथयति}} (\emph{kathayati}) — \textquotedblleft{}(he, she) tells\textquotedblright{}. A \textbf{\sk{कथा}}, \textquotedblleft{}a story\textquotedblright{}, is what is told.

\begin{grammarlens}[title={What you've built}]
The verb of a story.
\end{grammarlens}

\subsection*{Your turn}
\begin{itemize}
\raggedright
  \item \emph{Say it:} \emph{kathayati}
  \item \emph{Say it:} \emph{kathayati}, once more
  \item \emph{Say it:} say \emph{kathayati}
  \item \emph{Recall:} say the Sanskrit for she gives, then the Sanskrit for she sings, then say \emph{kathayati} again
\end{itemize}

\begin{tcolorbox}[breakable,colback=teal!4,colframe=teal!35!black,title={Before you move on}]
What does \sk{कथयति} mean? (\textquotedblleft{}Tells\textquotedblright{}.) Say it once more.
\end{tcolorbox}
\section[roditi]{\sk{रोदिति} (roditi) — cries}

\label{lesson:SA-C73-cries}



\begin{quote}
Before the new one: say the Sanskrit for she sings, then the Sanskrit for he tells.
\end{quote}

\subsection*{You'll want to know: \sk{रोदिति}}
\textbf{\sk{रोदिति}} (\emph{roditi}) — \textquotedblleft{}(he, she) cries\textquotedblright{}, the verb of \textbf{\sk{दुःखम्}}.

\begin{grammarlens}[title={What you've built}]
Sorrow, as a verb.
\end{grammarlens}

\subsection*{Your turn}
\begin{itemize}
\raggedright
  \item \emph{Say it:} \emph{roditi}
  \item \emph{Say it:} \emph{roditi}, once more
  \item \emph{Say it:} say \emph{bālaḥ roditi}
  \item \emph{Recall:} say the Sanskrit for she sings, then the Sanskrit for he tells, then say \emph{roditi} again
\end{itemize}

\begin{tcolorbox}[breakable,colback=teal!4,colframe=teal!35!black,title={Before you move on}]
What does \sk{रोदिति} mean? (\textquotedblleft{}Cries\textquotedblright{}.) Say it once more.
\end{tcolorbox}
\section[āhvayati]{\sk{आह्वयति} (āhvayati) — calls}

\label{lesson:SA-C73-calls}



\begin{quote}
Before the new one: say the Sanskrit for he tells, then the Sanskrit for the child cries.
\end{quote}

\subsection*{You'll want to know: \sk{आह्वयति}}
\textbf{\sk{आह्वयति}} (\emph{āhvayati}) — \textquotedblleft{}(he, she) calls (someone)\textquotedblright{}. \textbf{\sk{माता} \sk{आह्वयति}} — \textquotedblleft{}mother is calling\textquotedblright{}.

\begin{grammarlens}[title={What you've built}]
What you do across the courtyard.
\end{grammarlens}

\subsection*{Your turn}
\begin{itemize}
\raggedright
  \item \emph{Say it:} \emph{āhvayati}
  \item \emph{Say it:} \emph{āhvayati}, once more
  \item \emph{Say it:} say \emph{mātā āhvayati}
  \item \emph{Recall:} say the Sanskrit for he tells, then the Sanskrit for the child cries, then say \emph{āhvayati} again
\end{itemize}

\begin{tcolorbox}[breakable,colback=teal!4,colframe=teal!35!black,title={Before you move on}]
What does \sk{आह्वयति} mean? (\textquotedblleft{}Calls\textquotedblright{}.) Say it once more.
\end{tcolorbox}
\section[namati]{\sk{नमति} (namati) — bows}

\label{lesson:SA-C73-bows}



\begin{quote}
Before the new one: say the Sanskrit for the child cries, then the Sanskrit for she calls.
\end{quote}

\subsection*{You'll want to know: \sk{नमति}}
\textbf{\sk{नमति}} (\emph{namati}) — \textquotedblleft{}(he, she) bows\textquotedblright{}. It is the verb inside \textbf{\sk{नमस्ते}} and \textbf{\sk{प्रणामः}}.

\begin{grammarlens}[title={What you've built}]
The verb inside the greeting.
\end{grammarlens}

\subsection*{Your turn}
\begin{itemize}
\raggedright
  \item \emph{Say it:} \emph{namati}
  \item \emph{Say it:} \emph{namati}, once more
  \item \emph{Say it:} say \emph{namati}
  \item \emph{Recall:} say the Sanskrit for the child cries, then the Sanskrit for she calls, then say \emph{namati} again
\end{itemize}

\begin{tcolorbox}[breakable,colback=teal!4,colframe=teal!35!black,title={Before you move on}]
What does \sk{नमति} mean? (\textquotedblleft{}Bows\textquotedblright{}.) Say it once more.
\end{tcolorbox}
